% =========================================================
% 03_linear_combinations.tex
% Vectors and Linear Algebra
% =========================================================

% =========================================================
% SECTION DIVIDER
% =========================================================

\section{Linear Combinations and Span}

\begin{frame}{Linear Combinations}

A \textbf{linear combination} of vectors is formed by multiplying
vectors by scalars and adding the results.

For vectors $\vect{v}_1,\vect{v}_2,\ldots,\vect{v}_k$,

\[
\boxed{
\vect{w}
=
c_1\vect{v}_1
+
c_2\vect{v}_2
+\cdots+
c_k\vect{v}_k
}
\]

where $c_1,c_2,\ldots,c_k$ are scalars.

\vspace{0.4cm}

The scalars are called the \textbf{coefficients} of the linear combination.

\end{frame}


% =========================================================
% EXAMPLE
% =========================================================

\begin{frame}{Linear Combination: Example}

Let
\vspace{-0.1cm}
\[
\vect{u}
=
\begin{bmatrix}
1\\
2
\end{bmatrix},
\qquad
\vect{v}
=
\begin{bmatrix}
3\\
1
\end{bmatrix}.
\]

Consider
\vspace{-0.1cm}
\[
2\vect{u}-\vect{v}.
\]

Then
\vspace{-0.1cm}
\[
\begin{aligned}
2\vect{u}-\vect{v}
&=
2
\begin{bmatrix}
1\\
2
\end{bmatrix}
-
\begin{bmatrix}
3\\
1
\end{bmatrix}\\[4pt]
&=
\begin{bmatrix}
2\\
4
\end{bmatrix}
-
\begin{bmatrix}
3\\
1
\end{bmatrix}\\[4pt]
&=
\boxed{
\begin{bmatrix}
-1\\
3
\end{bmatrix}}
\end{aligned}
\]

\end{frame}


% =========================================================
% GEOMETRIC INTERPRETATION
% =========================================================

\begin{frame}{Linear Combinations: Geometric View}

Suppose

\[
\vect{w}=a\vect{u}+b\vect{v}.
\]

The coefficients $a$ and $b$ determine how much of each vector
contributes to $\vect{w}$.

\vspace{0.3cm}

\begin{center}

\begin{tikzpicture}[
    >=Latex,
    scale=0.75
]

\coordinate (O) at (0,0);
\coordinate (U) at (3,1);
\coordinate (V) at (1,2);
\coordinate (W) at (4,3);

\draw[->, very thick, PastelBlue]
    (O) -- (U)
    node[midway, below] {$\vect{u}$};

\draw[->, very thick, PastelLavender]
    (O) -- (V)
    node[midway, left] {$\vect{v}$};

\draw[dashed, gray] (U) -- (W);
\draw[dashed, gray] (V) -- (W);

\draw[->, very thick]
    (O) -- (W)
    node[midway, above] {$a\vect{u}+b\vect{v}$};

\end{tikzpicture}

\end{center}

\[
\boxed{\vect{w}=a\vect{u}+b\vect{v}}
\]

\end{frame}


% =========================================================
% SPAN
% =========================================================

\begin{frame}{Span of Vectors}

The \textbf{span} of a set of vectors is the set of all
possible linear combinations of those vectors.

For

\[
\vect{v}_1,\vect{v}_2,\ldots,\vect{v}_k,
\]

the span is

\[
\boxed{
\operatorname{span}
\{\vect{v}_1,\ldots,\vect{v}_k\}
=
\left\{
\sum_{i=1}^{k}c_i\vect{v}_i
\;:\;
c_i\in\mathbb{R}
\right\}
}
\]

\vspace{0.4cm}

The span describes the set of vectors that can be generated
from the given vectors.

\end{frame}


% =========================================================
% SPAN IN R2
% =========================================================

\begin{frame}{Span in $\mathbb{R}^2$}

Consider two vectors
\vspace{0.1cm}
\[
\vect{u}
=
\begin{bmatrix}
1\\
0
\end{bmatrix},
\qquad
\vect{v}
=
\begin{bmatrix}
0\\
1
\end{bmatrix}.
\]

Any vector
\vspace{0.1cm}
\[
\begin{bmatrix}
x\\
y
\end{bmatrix}
\in\mathbb{R}^2
\]

can be written as
\vspace{0.1cm}
\[
\begin{bmatrix}
x\\
y
\end{bmatrix}
=
x
\begin{bmatrix}
1\\
0
\end{bmatrix}
+
y
\begin{bmatrix}
0\\
1
\end{bmatrix}.
\]

Therefore,
\vspace{0.1cm}
\[
\boxed{
\operatorname{span}\{\vect{u},\vect{v}\}
=
\mathbb{R}^2
}
\]

\end{frame}


% =========================================================
% SPAN GEOMETRY
% =========================================================

\begin{frame}{Geometric Meaning of Span}

The span of vectors can have different dimensions.

\vspace{0.3cm}

\begin{columns}[T,totalwidth=\textwidth]

\begin{column}{0.31\textwidth}

\begin{block}{One Direction}

A non-zero vector spans a \textbf{line}.

\[
\operatorname{span}\{\vect{v}\}
\]

\end{block}

\end{column}

\begin{column}{0.31\textwidth}

\begin{block}{Two Directions}

Two independent vectors in $\mathbb{R}^2$ span a \textbf{plane}.

\[
\operatorname{span}\{\vect{u},\vect{v}\}
\]

\end{block}

\end{column}

\begin{column}{0.31\textwidth}

\begin{block}{Three Directions}

Three independent vectors in $\mathbb{R}^3$ span
three-dimensional space.

\[
\mathbb{R}^3
\]

\end{block}

\end{column}

\end{columns}

\vspace{0.5cm}

\begin{center}
The span describes the \textbf{space generated by the vectors}.
\end{center}

\end{frame}


% =========================================================
% LINEAR INDEPENDENCE
% =========================================================

\begin{frame}{Linearly Independent Vectors}

Vectors $\vect{v}_1,\ldots,\vect{v}_k$ are \textbf{linearly independent}
if the equation

\[
c_1\vect{v}_1+
c_2\vect{v}_2+
\cdots+
c_k\vect{v}_k
=
\mathbf{0}
\]

has only the trivial solution

\[
\boxed{
c_1=c_2=\cdots=c_k=0.
}
\]

\vspace{0.5cm}

In other words, none of the vectors can be written as a
linear combination of the others.

\end{frame}


% =========================================================
% LINEAR DEPENDENCE
% =========================================================

\begin{frame}{Linearly Dependent Vectors}

Vectors are \textbf{linearly dependent} if there exists a
non-trivial set of coefficients such that

\[
c_1\vect{v}_1+
c_2\vect{v}_2+
\cdots+
c_k\vect{v}_k
=
\mathbf{0},
\]

where at least one coefficient is non-zero.

\[
\boxed{
\text{Non-trivial solution}
\quad\Longrightarrow\quad
\text{Linear dependence}
}
\]

\vspace{0.4cm}

This means that at least one vector can be expressed as a
linear combination of the others.

\end{frame}


% =========================================================
% DEPENDENCE EXAMPLE
% =========================================================

\begin{frame}{Linear Dependence: Example}

Consider
\vspace{-0.1cm}
\[
\vect{u}
=
\begin{bmatrix}
1\\
2
\end{bmatrix},
\qquad
\vect{v}
=
\begin{bmatrix}
2\\
4
\end{bmatrix}.
\]

Notice that
\vspace{-0.1cm}
\[
\vect{v}=2\vect{u}.
\]

Therefore,
\vspace{-0.1cm}
\[
2\vect{u}-\vect{v}
=
\mathbf{0}.
\]

The coefficients
\vspace{-0.1cm}
\[
2\neq0,
\qquad
-1\neq0
\]

give a non-trivial solution.

Hence,

\[
\boxed{
\vect{u},\vect{v}
\text{ are linearly dependent.}
}
\]

\end{frame}


% =========================================================
% INDEPENDENCE EXAMPLE
% =========================================================

\begin{frame}{Linear Independence: Example}

Consider
\vspace{-0.1cm}
\[
\vect{u}
=
\begin{bmatrix}
1\\
0
\end{bmatrix},
\qquad
\vect{v}
=
\begin{bmatrix}
0\\
1
\end{bmatrix}.
\]

Suppose

\[
a\vect{u}+b\vect{v}
=
\mathbf{0}.
\]

Then
\vspace{-0.1cm}
\[
a
\begin{bmatrix}
1\\
0
\end{bmatrix}
+
b
\begin{bmatrix}
0\\
1
\end{bmatrix}
=
\begin{bmatrix}
0\\
0
\end{bmatrix}.
\]

Therefore,

\[
a=0,
\qquad
b=0.
\]

Hence,
\vspace{-0.1cm}
\[
\boxed{
\vect{u},\vect{v}
\text{ are linearly independent.}
}
\]

\end{frame}


% =========================================================
% LINEAR COMBINATION, SPAN, INDEPENDENCE
% =========================================================

\begin{frame}{How the Concepts Are Connected}

The ideas introduced so far are closely related:

\vspace{0.4cm}

\begin{center}

\begin{tikzpicture}[
    node distance=1.7cm,
    >=Latex
]

\node[
    draw,
    rounded corners,
    fill=LightBlue,
    minimum width=2.6cm,
    minimum height=0.8cm
] (vectors) {Vectors};

\node[
    draw,
    rounded corners,
    fill=LightBlue,
    minimum width=2.6cm,
    minimum height=0.8cm,
    right=of vectors
] (combination) {Linear Combination};

\node[
    draw,
    rounded corners,
    fill=LightLavender,
    minimum width=2.6cm,
    minimum height=0.8cm,
    right=of combination
] (span) {Span};

\node[
    draw,
    rounded corners,
    fill=LightLavender,
    minimum width=2.6cm,
    minimum height=0.8cm,
    below=of combination
] (independence) {Independence};

\draw[->, thick] (vectors) -- (combination);
\draw[->, thick] (combination) -- (span);
\draw[->, thick] (vectors) -- (independence);

\end{tikzpicture}

\end{center}

\vspace{0.4cm}

\begin{itemize}
    \item Linear combinations generate vectors.
    \item The span contains all vectors that can be generated.
    \item Linear independence describes whether the generators are redundant.
\end{itemize}

\end{frame}


% =========================================================
% LINEAR COMBINATION AND SYSTEMS
% =========================================================

\begin{frame}{Linear Combinations and Linear Systems}

Consider
\vspace{-0.1cm}
\[
A\vect{x}=\vect{b}.
\]

If
\vspace{-0.1cm}
\[
A=
\begin{bmatrix}
| & | & & |\\
\vect{a}_1 & \vect{a}_2 & \cdots & \vect{a}_n\\
| & | & & |
\end{bmatrix},
\]

then
\vspace{-0.1cm}
\[
A\vect{x}
=
x_1\vect{a}_1+
x_2\vect{a}_2+
\cdots+
x_n\vect{a}_n.
\]

Therefore,
\vspace{-0.1cm}
\[
\boxed{
A\vect{x}=\vect{b}
}
\]

asks whether $\vect{b}$ can be expressed as a linear combination
of the columns of $A$.

\vspace{0.1cm}

This connection will become important when we study
\textbf{column space and systems of linear equations}.

\end{frame}


% =========================================================
% MAXIMUM NUMBER OF INDEPENDENT VECTORS
% =========================================================

\begin{frame}{Independent Vectors and Dimension}

In $\mathbb{R}^n$, there can be at most $n$ linearly independent
vectors.

For example:

\[
\mathbb{R}^2
\quad\Rightarrow\quad
\text{at most 2 independent vectors}
\]

\[
\mathbb{R}^3
\quad\Rightarrow\quad
\text{at most 3 independent vectors}
\]

\[
\mathbb{R}^n
\quad\Rightarrow\quad
\text{at most }n\text{ independent vectors}.
\]

\vspace{0.5cm}

This observation leads naturally to the concepts of

\[
\boxed{
\text{Basis}
\qquad\text{and}\qquad
\text{Dimension}.
}
\]

\end{frame}